% Cuerpo del informe (sin preámbulo) para incorporar en otro documento LaTeX:
%   % Cuerpo del informe (sin preámbulo) para incorporar en otro documento LaTeX:
%   % Cuerpo del informe (sin preámbulo) para incorporar en otro documento LaTeX:
%   % Cuerpo del informe (sin preámbulo) para incorporar en otro documento LaTeX:
%   \input{docs/informe_proyecto_cuerpo}
% Requiere en el documento padre (según uso): babel spanish, booktabs, listings,
% enumitem, hyperref (opcionales pero recomendados).

\section{Introducción}
El proyecto académico \emph{Sistema de Autenticación y Autorización Centralizado
(Master Gateway)} responde a la necesidad de una puerta maestra de seguridad
para arquitecturas de microservicios. En lugar de que cada servicio gestione
usuarios y credenciales, el Master autentica usuarios, obliga a seleccionar un
rol activo, emite JWT con privilegio mínimo, expone validación interna Zero Trust
para microservicios hijos y administra usuarios, roles, módulos, menús y permisos
con soft delete y auditoría.

Principios rectores: Zero Trust, Least Privilege, Shift-Left Security,
autenticación stateless y soft delete.

\section{Objetivos}
\begin{enumerate}
  \item Implementar el flujo login $\rightarrow$ TempToken $\rightarrow$ selección de rol
        $\rightarrow$ AccessToken + RefreshToken rotativo.
  \item Modelar entidades con auditoría y estado ACTIVO/INACTIVO.
  \item Generar menú dinámico con CTE recursiva filtrado por rol JWT.
  \item Integrar un microservicio hijo (Ventas) bajo Zero Trust.
  \item Entregar SPA con rutas protegidas y navegación dinámica.
  \item Automatizar calidad y despliegue (GitHub Actions, SonarCloud, SAST/ML, Railway, Telegram).
\end{enumerate}

\section{Arquitectura del sistema}
\subsection{Vista lógica}
\begin{enumerate}
  \item \textbf{Master Gateway}: API Node.js/TypeScript (auth, RBAC, CRUD, validate-token).
  \item \textbf{PostgreSQL} vía Prisma ORM.
  \item \textbf{SPA React + Vite}: login, selector de rol, administración y Ventas.
  \item \textbf{Microservicio Ventas}: sin DB de usuarios; valida tokens en el Master.
\end{enumerate}

\subsection{Flujo Zero Trust}
El frontend envía el AccessToken a Ventas; Ventas llama a
\texttt{POST /api/internals/validate-token} con \texttt{x-internal-api-key}.
Solo si el Master confirma validez y permisos se ejecuta la lógica de negocio.

\subsection{Stack}
Master: Express, TypeScript, Prisma, Argon2, Jose (JWT), Helmet, Zod.
Frontend: React 19, Vite, React Router.
CI/CD: GitHub Actions, SonarCloud, SAST/ML, Railway CLI, Telegram.

\section{Modelo de datos}
Entidades: User, Role, UserRole, Module, RoleModule, Menu (Adjacency List),
RoleMenu, Permission, RolePermission, RefreshToken, TokenRevocation, AuditLog.
Menú: solo hojas con URL, anticiclos, exclusión de inactivos.

\section{Autenticación y autorización}
Login emite TempToken + roles; la SPA fuerza select-role; el AccessToken incluye
únicamente el rol seleccionado y sus permisos. Refresh rotativo y logout con
revocación por \texttt{jti}. Endpoint interno
\texttt{/api/internals/validate-token} responde
\texttt{active}, \texttt{userId}, \texttt{roleId}, \texttt{roleName}, \texttt{permissions}.

\section{Frontend}
SPA con sessionStorage, interceptor Bearer/refresh, menú dinámico, administración
de usuarios/roles, pantallas 403/sesión/token expirado y cancelación del selector
de rol para cambiar de usuario. La UI no decide autorización final.

\section{Microservicio Ventas}
Endpoints con permisos \texttt{VENTAS\_READ}/\texttt{VENTAS\_CREATE}; retries ante
cold start PaaS. Roles demo: ADMIN, VENDEDOR, AUDITOR.

\section{Seguridad}
Argon2, rate limiting, Helmet, Zod, soft delete, secretos en entorno, mensajes
de login genéricos. JWT HS256 + validate-token (RS256 opcional no bloqueante).

\section{DevSecOps}
Ramas \texttt{feature $\rightarrow$ dev $\rightarrow$ test $\rightarrow$ main}.
Pipeline en merge a main: build/tests, Sonar Quality Gate obligatorio, SAST/ML,
deploy Railway, Telegram. Despliegue recomendado: Railway (alternativa Render).

\section{Pruebas}
Vitest, smoke E2E (\texttt{npm run smoke}), checklist en
\texttt{docs/COMPLIANCE\_CHECKLIST.md}.

\section{Conclusiones}
Se cumple el alcance académico Full-Stack seguro (Master + SPA + Ventas + DevSecOps).
Pendiente operativo: secrets GitHub, branch protection, SonarCloud y PaaS.
Mejora opcional: JWT asimétrico.

% Requiere en el documento padre (según uso): babel spanish, booktabs, listings,
% enumitem, hyperref (opcionales pero recomendados).

\section{Introducción}
El proyecto académico \emph{Sistema de Autenticación y Autorización Centralizado
(Master Gateway)} responde a la necesidad de una puerta maestra de seguridad
para arquitecturas de microservicios. En lugar de que cada servicio gestione
usuarios y credenciales, el Master autentica usuarios, obliga a seleccionar un
rol activo, emite JWT con privilegio mínimo, expone validación interna Zero Trust
para microservicios hijos y administra usuarios, roles, módulos, menús y permisos
con soft delete y auditoría.

Principios rectores: Zero Trust, Least Privilege, Shift-Left Security,
autenticación stateless y soft delete.

\section{Objetivos}
\begin{enumerate}
  \item Implementar el flujo login $\rightarrow$ TempToken $\rightarrow$ selección de rol
        $\rightarrow$ AccessToken + RefreshToken rotativo.
  \item Modelar entidades con auditoría y estado ACTIVO/INACTIVO.
  \item Generar menú dinámico con CTE recursiva filtrado por rol JWT.
  \item Integrar un microservicio hijo (Ventas) bajo Zero Trust.
  \item Entregar SPA con rutas protegidas y navegación dinámica.
  \item Automatizar calidad y despliegue (GitHub Actions, SonarCloud, SAST/ML, Railway, Telegram).
\end{enumerate}

\section{Arquitectura del sistema}
\subsection{Vista lógica}
\begin{enumerate}
  \item \textbf{Master Gateway}: API Node.js/TypeScript (auth, RBAC, CRUD, validate-token).
  \item \textbf{PostgreSQL} vía Prisma ORM.
  \item \textbf{SPA React + Vite}: login, selector de rol, administración y Ventas.
  \item \textbf{Microservicio Ventas}: sin DB de usuarios; valida tokens en el Master.
\end{enumerate}

\subsection{Flujo Zero Trust}
El frontend envía el AccessToken a Ventas; Ventas llama a
\texttt{POST /api/internals/validate-token} con \texttt{x-internal-api-key}.
Solo si el Master confirma validez y permisos se ejecuta la lógica de negocio.

\subsection{Stack}
Master: Express, TypeScript, Prisma, Argon2, Jose (JWT), Helmet, Zod.
Frontend: React 19, Vite, React Router.
CI/CD: GitHub Actions, SonarCloud, SAST/ML, Railway CLI, Telegram.

\section{Modelo de datos}
Entidades: User, Role, UserRole, Module, RoleModule, Menu (Adjacency List),
RoleMenu, Permission, RolePermission, RefreshToken, TokenRevocation, AuditLog.
Menú: solo hojas con URL, anticiclos, exclusión de inactivos.

\section{Autenticación y autorización}
Login emite TempToken + roles; la SPA fuerza select-role; el AccessToken incluye
únicamente el rol seleccionado y sus permisos. Refresh rotativo y logout con
revocación por \texttt{jti}. Endpoint interno
\texttt{/api/internals/validate-token} responde
\texttt{active}, \texttt{userId}, \texttt{roleId}, \texttt{roleName}, \texttt{permissions}.

\section{Frontend}
SPA con sessionStorage, interceptor Bearer/refresh, menú dinámico, administración
de usuarios/roles, pantallas 403/sesión/token expirado y cancelación del selector
de rol para cambiar de usuario. La UI no decide autorización final.

\section{Microservicio Ventas}
Endpoints con permisos \texttt{VENTAS\_READ}/\texttt{VENTAS\_CREATE}; retries ante
cold start PaaS. Roles demo: ADMIN, VENDEDOR, AUDITOR.

\section{Seguridad}
Argon2, rate limiting, Helmet, Zod, soft delete, secretos en entorno, mensajes
de login genéricos. JWT HS256 + validate-token (RS256 opcional no bloqueante).

\section{DevSecOps}
Ramas \texttt{feature $\rightarrow$ dev $\rightarrow$ test $\rightarrow$ main}.
Pipeline en merge a main: build/tests, Sonar Quality Gate obligatorio, SAST/ML,
deploy Railway, Telegram. Despliegue recomendado: Railway (alternativa Render).

\section{Pruebas}
Vitest, smoke E2E (\texttt{npm run smoke}), checklist en
\texttt{docs/COMPLIANCE\_CHECKLIST.md}.

\section{Conclusiones}
Se cumple el alcance académico Full-Stack seguro (Master + SPA + Ventas + DevSecOps).
Pendiente operativo: secrets GitHub, branch protection, SonarCloud y PaaS.
Mejora opcional: JWT asimétrico.

% Requiere en el documento padre (según uso): babel spanish, booktabs, listings,
% enumitem, hyperref (opcionales pero recomendados).

\section{Introducción}
El proyecto académico \emph{Sistema de Autenticación y Autorización Centralizado
(Master Gateway)} responde a la necesidad de una puerta maestra de seguridad
para arquitecturas de microservicios. En lugar de que cada servicio gestione
usuarios y credenciales, el Master autentica usuarios, obliga a seleccionar un
rol activo, emite JWT con privilegio mínimo, expone validación interna Zero Trust
para microservicios hijos y administra usuarios, roles, módulos, menús y permisos
con soft delete y auditoría.

Principios rectores: Zero Trust, Least Privilege, Shift-Left Security,
autenticación stateless y soft delete.

\section{Objetivos}
\begin{enumerate}
  \item Implementar el flujo login $\rightarrow$ TempToken $\rightarrow$ selección de rol
        $\rightarrow$ AccessToken + RefreshToken rotativo.
  \item Modelar entidades con auditoría y estado ACTIVO/INACTIVO.
  \item Generar menú dinámico con CTE recursiva filtrado por rol JWT.
  \item Integrar un microservicio hijo (Ventas) bajo Zero Trust.
  \item Entregar SPA con rutas protegidas y navegación dinámica.
  \item Automatizar calidad y despliegue (GitHub Actions, SonarCloud, SAST/ML, Railway, Telegram).
\end{enumerate}

\section{Arquitectura del sistema}
\subsection{Vista lógica}
\begin{enumerate}
  \item \textbf{Master Gateway}: API Node.js/TypeScript (auth, RBAC, CRUD, validate-token).
  \item \textbf{PostgreSQL} vía Prisma ORM.
  \item \textbf{SPA React + Vite}: login, selector de rol, administración y Ventas.
  \item \textbf{Microservicio Ventas}: sin DB de usuarios; valida tokens en el Master.
\end{enumerate}

\subsection{Flujo Zero Trust}
El frontend envía el AccessToken a Ventas; Ventas llama a
\texttt{POST /api/internals/validate-token} con \texttt{x-internal-api-key}.
Solo si el Master confirma validez y permisos se ejecuta la lógica de negocio.

\subsection{Stack}
Master: Express, TypeScript, Prisma, Argon2, Jose (JWT), Helmet, Zod.
Frontend: React 19, Vite, React Router.
CI/CD: GitHub Actions, SonarCloud, SAST/ML, Railway CLI, Telegram.

\section{Modelo de datos}
Entidades: User, Role, UserRole, Module, RoleModule, Menu (Adjacency List),
RoleMenu, Permission, RolePermission, RefreshToken, TokenRevocation, AuditLog.
Menú: solo hojas con URL, anticiclos, exclusión de inactivos.

\section{Autenticación y autorización}
Login emite TempToken + roles; la SPA fuerza select-role; el AccessToken incluye
únicamente el rol seleccionado y sus permisos. Refresh rotativo y logout con
revocación por \texttt{jti}. Endpoint interno
\texttt{/api/internals/validate-token} responde
\texttt{active}, \texttt{userId}, \texttt{roleId}, \texttt{roleName}, \texttt{permissions}.

\section{Frontend}
SPA con sessionStorage, interceptor Bearer/refresh, menú dinámico, administración
de usuarios/roles, pantallas 403/sesión/token expirado y cancelación del selector
de rol para cambiar de usuario. La UI no decide autorización final.

\section{Microservicio Ventas}
Endpoints con permisos \texttt{VENTAS\_READ}/\texttt{VENTAS\_CREATE}; retries ante
cold start PaaS. Roles demo: ADMIN, VENDEDOR, AUDITOR.

\section{Seguridad}
Argon2, rate limiting, Helmet, Zod, soft delete, secretos en entorno, mensajes
de login genéricos. JWT HS256 + validate-token (RS256 opcional no bloqueante).

\section{DevSecOps}
Ramas \texttt{feature $\rightarrow$ dev $\rightarrow$ test $\rightarrow$ main}.
Pipeline en merge a main: build/tests, Sonar Quality Gate obligatorio, SAST/ML,
deploy Railway, Telegram. Despliegue recomendado: Railway (alternativa Render).

\section{Pruebas}
Vitest, smoke E2E (\texttt{npm run smoke}), checklist en
\texttt{docs/COMPLIANCE\_CHECKLIST.md}.

\section{Conclusiones}
Se cumple el alcance académico Full-Stack seguro (Master + SPA + Ventas + DevSecOps).
Pendiente operativo: secrets GitHub, branch protection, SonarCloud y PaaS.
Mejora opcional: JWT asimétrico.

% Requiere en el documento padre (según uso): babel spanish, booktabs, listings,
% enumitem, hyperref (opcionales pero recomendados).

\section{Introducción}
El proyecto académico \emph{Sistema de Autenticación y Autorización Centralizado
(Master Gateway)} responde a la necesidad de una puerta maestra de seguridad
para arquitecturas de microservicios. En lugar de que cada servicio gestione
usuarios y credenciales, el Master autentica usuarios, obliga a seleccionar un
rol activo, emite JWT con privilegio mínimo, expone validación interna Zero Trust
para microservicios hijos y administra usuarios, roles, módulos, menús y permisos
con soft delete y auditoría.

Principios rectores: Zero Trust, Least Privilege, Shift-Left Security,
autenticación stateless y soft delete.

\section{Objetivos}
\begin{enumerate}
  \item Implementar el flujo login $\rightarrow$ TempToken $\rightarrow$ selección de rol
        $\rightarrow$ AccessToken + RefreshToken rotativo.
  \item Modelar entidades con auditoría y estado ACTIVO/INACTIVO.
  \item Generar menú dinámico con CTE recursiva filtrado por rol JWT.
  \item Integrar un microservicio hijo (Ventas) bajo Zero Trust.
  \item Entregar SPA con rutas protegidas y navegación dinámica.
  \item Automatizar calidad y despliegue (GitHub Actions, SonarCloud, SAST/ML, Railway, Telegram).
\end{enumerate}

\section{Arquitectura del sistema}
\subsection{Vista lógica}
\begin{enumerate}
  \item \textbf{Master Gateway}: API Node.js/TypeScript (auth, RBAC, CRUD, validate-token).
  \item \textbf{PostgreSQL} vía Prisma ORM.
  \item \textbf{SPA React + Vite}: login, selector de rol, administración y Ventas.
  \item \textbf{Microservicio Ventas}: sin DB de usuarios; valida tokens en el Master.
\end{enumerate}

\subsection{Flujo Zero Trust}
El frontend envía el AccessToken a Ventas; Ventas llama a
\texttt{POST /api/internals/validate-token} con \texttt{x-internal-api-key}.
Solo si el Master confirma validez y permisos se ejecuta la lógica de negocio.

\subsection{Stack}
Master: Express, TypeScript, Prisma, Argon2, Jose (JWT), Helmet, Zod.
Frontend: React 19, Vite, React Router.
CI/CD: GitHub Actions, SonarCloud, SAST/ML, Railway CLI, Telegram.

\section{Modelo de datos}
Entidades: User, Role, UserRole, Module, RoleModule, Menu (Adjacency List),
RoleMenu, Permission, RolePermission, RefreshToken, TokenRevocation, AuditLog.
Menú: solo hojas con URL, anticiclos, exclusión de inactivos.

\section{Autenticación y autorización}
Login emite TempToken + roles; la SPA fuerza select-role; el AccessToken incluye
únicamente el rol seleccionado y sus permisos. Refresh rotativo y logout con
revocación por \texttt{jti}. Endpoint interno
\texttt{/api/internals/validate-token} responde
\texttt{active}, \texttt{userId}, \texttt{roleId}, \texttt{roleName}, \texttt{permissions}.

\section{Frontend}
SPA con sessionStorage, interceptor Bearer/refresh, menú dinámico, administración
de usuarios/roles, pantallas 403/sesión/token expirado y cancelación del selector
de rol para cambiar de usuario. La UI no decide autorización final.

\section{Microservicio Ventas}
Endpoints con permisos \texttt{VENTAS\_READ}/\texttt{VENTAS\_CREATE}; retries ante
cold start PaaS. Roles demo: ADMIN, VENDEDOR, AUDITOR.

\section{Seguridad}
Argon2, rate limiting, Helmet, Zod, soft delete, secretos en entorno, mensajes
de login genéricos. JWT HS256 + validate-token (RS256 opcional no bloqueante).

\section{DevSecOps}
Ramas \texttt{feature $\rightarrow$ dev $\rightarrow$ test $\rightarrow$ main}.
Pipeline en merge a main: build/tests, Sonar Quality Gate obligatorio, SAST/ML,
deploy Railway, Telegram. Despliegue recomendado: Railway (alternativa Render).

\section{Pruebas}
Vitest, smoke E2E (\texttt{npm run smoke}), checklist en
\texttt{docs/COMPLIANCE\_CHECKLIST.md}.

\section{Conclusiones}
Se cumple el alcance académico Full-Stack seguro (Master + SPA + Ventas + DevSecOps).
Pendiente operativo: secrets GitHub, branch protection, SonarCloud y PaaS.
Mejora opcional: JWT asimétrico.
